% =====================================================================
%  CHAPTER 3 — ELECTROMAGNETIC MODELLING OF THE RECEIVING ELEMENT
% =====================================================================
%  Compilation notes (pdflatex / VS Code, LaTeX Workshop):
%    Required packages (declare in your main preamble if not already):
%      \usepackage{siunitx}            % \SI, \si, \ohm, \mega\hertz, ...
%      \usepackage{booktabs}           % \toprule \midrule \bottomrule
%      \usepackage{enumitem}           % [leftmargin=...]
%      \usepackage{graphicx}           % \includegraphics
%      \usepackage{float}              % [H] float placement
%      \usepackage{amsmath,amssymb}    % equation, align, \mathrm, \approx
%      \usepackage[round]{natbib}      % \citep / \citet  (or biblatex; see notes)
%    Recommended siunitx setup (preamble):
%      \sisetup{detect-all, per-mode=symbol}
%    Bibliography keys used here (define them later in your .bib):
%      balanis2016, kraus1988, stutzman2013, pozar2012, harrington1993,
%      nec2, yagi1928, gupta1983, tms2017, mclean1985,
%      wild1963, benz2005, christiansen1985, sws_rims_metadata, rstn_overview
%  ---------------------------------------------------------------------
%  Figure files expected in your graphics path (placeholders — replace):
%      fig_yagi_geometry      fig_folded_dipole      fig_field_setup
%      fig_nec_model          fig_pattern_eplane     fig_pattern_3d
%      fig_coupling_curve     fig_primary_beam
% =====================================================================

\chapter{ELECTROMAGNETIC MODELLING OF THE RECEIVING ELEMENT}
\label{ch:antenna}

It is worth noting at the outset that low-cost interferometer projects such as
DLITE (Chapter~\ref{ch:intro}) favour LWA-style crossed-dipole antennas, which
are better suited to DLITE's much lower \SIrange{30}{40}{\mega\hertz} operating
band; the Yagi-Uda antenna adopted here is not an untested substitute chosen
for lack of an alternative. Yagi-class antennas have a long and independently
documented track record as a low-cost, high-reliability receiving element,
in amateur radio and in commercial UHF television reception alike
\citep{yagi1928,balanis2016,kraus1988}, and this chapter goes on to show that
the same combination of low cost and dependable performance carries over
directly to the \SI{610}{\mega\hertz} solar case: the specific unit used in
this thesis is validated below against three independent methods
(Section~\ref{sec:single-element-validation}) and found to be well matched,
with negligible mutual coupling at the operating baselines
(Section~\ref{sec:coupling-rationale}).

The accuracy of the entire interferometric simulation depends on having a
reliable model of the receiving element, and this chapter builds that model
as a chain of evidence rather than a survey of antenna theory: a real
antenna is characterised for gain and beam shape, its impedance is checked
for a stable match, and finally its behaviour when placed next to an
identical neighbour is checked against the independent-sensor assumption
that Chapter~\ref{ch:theory} relies on. Only once every link in that chain
holds is the antenna model handed to the simulation pipeline of
Chapter~\ref{ch:pipeline}. This thesis is, at its core, a study of an
\emph{antenna system}: the solar radio burst is the observation target, but
the main contribution is the careful characterisation of the receiving array
and the demonstration that its electromagnetic behaviour supports the
assumptions made in the signal-processing pipeline.

% =====================================================================
\section{Why the Receiving Element Must Be Modelled}
\label{sec:why-model}

In an aperture-synthesis array, the single-element radiation pattern is the
primary beam of the instrument (Section~\ref{subsec:antenna-gain-beam}). It
sets the field of view within which the Sun can be observed, and it controls
how strongly off-axis terrestrial interference enters the measurement.
Electromagnetic coupling between neighbouring elements also needs to be
verified: if coupling is significant, the array cannot be treated as a set
of independent sensors-the core assumption of the van~Cittert-Zernike
imaging framework developed in Chapter~\ref{ch:theory}-and a more complex
calibration would be required. Concretely, before any solar data can be
trusted, this chapter must answer four questions about the physical antenna:
\begin{enumerate}[leftmargin=1.8em]
  \item \textbf{Reception:} does the antenna actually receive a usable
        signal at \SI{610}{\mega\hertz}?
  \item \textbf{Beam:} what is its real radiation pattern, so that the
        correct primary beam-not an idealised Gaussian-can be used in the
        simulation?
  \item \textbf{Matching:} is its impedance stable and well matched enough
        that the received power computed by the Friis relation
        (Eq.~\eqref{eq:friis}) is a physically meaningful quantity, rather
        than being dominated by reflected power lost back down the feed?
  \item \textbf{Independence:} when several such antennas are placed close
        together to form an array, does mutual coupling break the
        independent-sensor assumption on which the imaging theory of
        Chapter~\ref{ch:theory} depends?
\end{enumerate}
The remainder of this chapter answers these four questions in order, in one
continuous chain: the receiving element is introduced and shown to be
suitable in principle (Sections~\ref{sec:yagi-selection}-\ref{sec:matching});
it is placed against a professional reference system and a field trial to
motivate why a controlled model is necessary at all
(Section~\ref{sec:reference-systems}); it is then modelled numerically
(Section~\ref{sec:simulation}) and validated against direct measurement
(Section~\ref{sec:single-element-validation}); its behaviour as an array
element is checked against the independent-sensor assumption
(Section~\ref{sec:coupling-rationale}); and finally its measured beam is
exported in the form actually consumed by the simulation
(Section~\ref{sec:primary-beam-extraction}).

% =====================================================================
\section{Yagi-Uda Antenna as the Selected Low-Cost Receiving Element}
\label{sec:yagi-selection}

The Yagi-Uda antenna remains one of the most widely used directional antennas
in the Ultra-High-Frequency (UHF) band. First described in the English-language
literature by \citet{yagi1928}, building on the work of Shintaro Uda, the design
achieves an attractive balance between high directivity, mechanical robustness,
and very low fabrication cost \citep{balanis2016,kraus1988}. These properties
are precisely what a low-budget, educational solar radio instrument requires,
and they are the reason a commercial Yagi forms the receiving element of this
project.

% ---------------------------------------------------------------------
\subsection{Operating Principle and Parasitic Excitation}
\label{subsec:parasitic}

Unlike a phased array, in which every element is fed with an individually
controlled voltage and phase, the Yagi-Uda array operates by
\emph{parasitic excitation}: it comprises a single active \emph{driven
element} together with several passive (parasitic) elements whose currents
are induced purely by electromagnetic coupling with their neighbours. This
is the same physical mechanism, generalised to many elements, that reappears
in Section~\ref{sec:coupling-rationale} between two separate antennas of the
array: a current flowing on one conductor induces a voltage on every nearby
conductor, and the strength of that induced voltage, for a given current, is
called the \emph{mutual impedance}.

For an array of $N+1$ elements, indexing the reflector as $-1$, the driven
element as $0$, and the directors as $1,2,\dots,N$, the terminal voltage on
element $m$ is the sum of the voltage its own current would produce if it
were isolated, plus the voltage induced on it by the current on every other
element:
\begin{equation}
    V_m = \sum_{n} Z_{mn}\, I_n ,
    \label{eq:mutual-matrix}
\end{equation}
where:
\begin{itemize}[leftmargin=1.5em,itemsep=1pt,topsep=2pt]
  \item $V_m$ is the terminal voltage of element $m$-the voltage that would
        be measured across its (real or virtual) feed gap;
  \item $I_n$ is the current on element $n$;
  \item $Z_{mn}$, for $m\neq n$, is the \emph{mutual impedance} between
        elements $m$ and $n$-the voltage induced on element $m$ per unit
        current on element $n$;
  \item $Z_{mm}$ is the \emph{self-impedance} of element $m$-the impedance
        it would present if every other element carried zero current
        \citep{balanis2016}.
\end{itemize}
Only the driven element is fed, so $V_0\neq 0$ while every parasitic element is
short-circuited at its (virtual) terminals:
\begin{equation}
    0 = \sum_{n} Z_{kn}\, I_n \qquad (k \neq 0).
    \label{eq:parasitic-zero}
\end{equation}
Equation~\eqref{eq:parasitic-zero} \citep{balanis2016,kraus1988} shows that the parasitic currents are not
free parameters: they are fixed by the geometry through the off-diagonal mutual
impedances. The design freedom available to the engineer is therefore the
choice of element lengths and spacings, which set the magnitude and-crucially
-the \emph{phase} of each induced current.

The directional behaviour follows directly from these induced phases
\citep{kraus1988,stutzman2013}:
\begin{itemize}[leftmargin=1.6em]
    \item \textbf{Reflector.} Positioned behind the driven element and made
    slightly longer than the resonant half-wavelength ($\approx\!5\%$ longer),
    the reflector presents an \emph{inductive} reactance. The induced current
    consequently lags the driving field, and the superposition of the
    re-radiated field with the primary field is constructive toward the front
    and destructive toward the rear. A single reflector captures most of the
    available rearward suppression; additional reflectors yield diminishing
    returns, which is why practical Yagis use one reflector (or, as here, one
    reflecting structure) but many directors.
    \item \textbf{Directors.} Positioned in front of the driven element and made
    $\approx\!5\%$ \emph{shorter} than resonance, the directors are
    \emph{capacitive}. Their induced currents lead the driving field, so the
    composite wave is progressively guided forward, sharpening the main lobe
    into an end-fire beam along the boom axis \citep{balanis2016}.
\end{itemize}
The net effect is a travelling-wave-like structure that concentrates radiated
(or, by reciprocity, received) power along the boom.

% ---------------------------------------------------------------------
\subsection{Directivity Build-Up and the Gain-Bandwidth Trade-off}
\label{subsec:directivity}

The directivity of a Yagi-Uda array increases monotonically, though
sub-linearly, with the number of directors. A useful qualitative progression,
consistent with both classical treatments and the present design, is that an
isolated half-wave dipole provides about \SI{2.15}{dBi}; adding a reflector
raises this to roughly \SI{7}{dBi}; the first director takes it close to
\SI{9.7}{dBi}; and a well-tuned six-director structure of the type used here
reaches the \SIrange{9.5}{13}{dBi} range depending on boom length and tuning
\citep{balanis2016,stutzman2013}. Each added director narrows the main lobe,
which improves angular resolution and-important for a ground-based solar
instrument-reduces the half-power beamwidth through which horizon-level
terrestrial interference can enter.

This gain, however, comes at the cost of bandwidth and tuning sensitivity. As
more directors are added, the array impedance becomes more reactive and the
input match deteriorates away from the design frequency. Pushing for too much
gain can reduce the usable bandwidth and even lower the peak gain by several
decibels if the element phasing moves away from the optimum \citep{stutzman2013}. For the solar application this trade-off is benign: the
instantaneous observing bandwidth is narrow (a single \SI{1}{\mega\hertz}
channel centred at \SI{610}{\mega\hertz}), so a moderately high-gain,
moderately broadband design is entirely appropriate, and the bandwidth inherited
from the original DVB-T2 service is far wider than the observation requires.

% ---------------------------------------------------------------------
\subsection{The Selected Physical Antenna: HKD~G11KD-T2}
\label{subsec:hkd-spec}

The physical receiving element adopted for this project is the commercial
HKD~G11KD-T2, a high-gain UHF Yagi-Uda antenna originally marketed for DVB-T2
digital-television reception. As a commercial-off-the-shelf (COTS) component it
combines mature RF engineering with rugged, weather-resistant aluminium
construction, and it is available at a small fraction of the cost of a
purpose-built radio-astronomy feed-the same low-element-cost argument
developed for the LWA-style antennas of DLITE in Chapter~\ref{ch:intro}.

Because the manufacturer publishes only marketing specifications, the geometry
used for simulation was established by reconciling three independent sources:
the manufacturer drawing, photographs of the deployed unit
(Figure~\ref{fig:hkd-photo}), and direct physical measurement of the antenna
with callipers. The resulting model, summarised in Table~\ref{tab:hkd-geom},
comprises an eleven-element structure on an aluminium boom of overall length
$\approx\!\SI{550}{\milli\metre}$. This same geometry is the input to the
4NEC2 antenna simulation software described in Section~\ref{sec:simulation}.

\begin{table}[H]
  \centering
  \caption[HKD G11KD-T2 antenna geometry used as electromagnetic model basis]{HKD G11KD-T2 antenna geometry used as electromagnetic model basis.\\[4pt]
           Dimensions were established by reconciling the manufacturer drawing,
           photographs of the deployed unit at the Xuan~Mai test site, and direct
           calliper measurements. The eleven-element Yagi sits on an aluminium
           boom of overall length approximately 550~mm; spacings are measured
           along the boom from the driven element (director~1 is the element
           closest to the feed).}
  \label{tab:hkd-geom}
  \renewcommand{\arraystretch}{1.25}
  \begin{tabularx}{\textwidth}{l >{\raggedright\arraybackslash}X}
    \toprule
    \textbf{Element} & \textbf{Dimension} \\
    \midrule
    Reflector (four-bar U-frame) & outer bars at $z=\pm\SI{40}{\milli\metre}$, inner bars at $z=\pm\SI{10}{\milli\metre}$ \\
    Driven element (folded dipole) & length $\approx\SI{210}{\milli\metre}$, fold height $\SI{44}{\milli\metre}$ \\
    Director lengths               & $\approx\SI{210}{\milli\metre}$ (each) \\
    Director spacings              & $82,\,40,\,50,\,50,\,66,\,69~\si{\milli\metre}$ \\
    Boom (square profile)          & $15\times15~\si{\milli\metre}$, length $\approx\SI{550}{\milli\metre}$ \\
    Element material               & aluminium (hollow tube / solid rod) \\
    \bottomrule
  \end{tabularx}
\end{table}

\begin{figure}[H]
  \centering
  \includegraphics[width=\textwidth]{photo_hkd_antenna}
  \caption[Physical HKD~G11KD-T2 Yagi-Uda antenna unit]{Physical HKD~G11KD-T2 Yagi-Uda antenna unit.\\[4pt]
  The eleven-element structure is shown with the four-bar U-frame reflector at the
  rear and the sequence of directors extending forward along the aluminium boom.
  The folded-dipole driven element sits immediately in front of the reflector frame;
  the mast-head LNA mounts directly at the feed terminals behind the driven element.
  Calliper measurements of this unit provided the dimensional inputs for
  Table~\ref{tab:hkd-geom}.}
  \label{fig:hkd-photo}
\end{figure}

The four-bar U-frame reflector is the distinguishing geometric feature of this
antenna. Rather than a single straight rod, the reflector forms a shallow
three-dimensional frame: this widens the effective reflecting aperture, raises
the front-to-back ratio, and in the radio-astronomy context suppresses pickup of
warm-ground thermal noise entering through the rear lobe.

% ---------------------------------------------------------------------
\subsection{Suitability for \texorpdfstring{\SI{610}{\mega\hertz}}{610 MHz}
Solar Observation}
\label{subsec:suitability}

At \SI{610}{\mega\hertz} the free-space wavelength is
\begin{equation}
    \lambda = \frac{c}{f} = \frac{2.998\times10^{8}\,\si{\metre\per\second}}
                                 {610\times10^{6}\,\si{\hertz}}
            \approx \SI{491.5}{\milli\metre},
    \label{eq:wavelength}
\end{equation}
so a multi-element Yagi remains physically compact (boom length of order one
metre) and can be hand-built and tripod-mounted-a decisive advantage for a
low-cost, portable instrument \citep{balanis2016}. The Sun is a weak source at
this frequency during quiet conditions (Chapter~\ref{ch:theory}), and any
target burst must be separated from a strong terrestrial noise background. A
directional element with moderate gain ($\sim\!\SI{9.5}{dBi}$) and a high
front-to-back ratio is therefore valuable: it concentrates sensitivity toward
the solar transit path while suppressing pickup from the horizon and from
behind the array \citep{tms2017}.

There is, however, an important double-edged aspect to using a
commercial DVB-T2 antenna for this purpose, and it is worth stating explicitly
because it shapes the field results in Section~\ref{sec:reference-systems}. The very
features that make the antenna attractive-that it is mass-produced, inexpensive,
and resonant across the \SIrange{470}{862}{\mega\hertz} UHF television band-
also make it maximally sensitive to terrestrial digital-television broadcasts
that occupy exactly that band. In a populated area the \SI{610}{\mega\hertz}
observing channel sits in the middle of an active broadcast allocation, so the
antenna efficiently collects precisely the interference one most wishes to
avoid. This tension between availability and radio-frequency cleanliness is a
recurring theme of the field campaign and one of the principal motivations for
retreating to a controlled numerical model.

% =====================================================================
\section{Impedance, Folded Dipole, and 75-\texorpdfstring{$\Omega$}{Ohm} Feed Matching}
\label{sec:matching}

Before the antenna can be simulated or measured, it is worth being precise
about which impedance is which, because three different quantities appear in
this chapter and they are easy to conflate.
\begin{itemize}[leftmargin=1.5em,itemsep=1pt,topsep=2pt]
  \item $Z_d$ is the input impedance (equivalently, the radiation resistance)
        of a single, isolated half-wave dipole under idealised free-space
        conditions. The commonly used approximate value is
        $Z_d\approx\SI{73}{\ohm}$.
  \item $Z_f$ is the input impedance of a \emph{folded} dipole: two closely
        spaced parallel conductors joined at both ends, fed by breaking only
        one conductor at its centre.
  \item $Z_\mathrm{in}$ is the actual input impedance measured, or simulated,
        at the physical feed port of the \emph{complete} antenna-after the
        reflector, all six directors, and their mutual loading are included.
        $Z_\mathrm{in}$ is not the same quantity as $Z_f$, and it is
        $Z_\mathrm{in}$, not $Z_f$, that is reported for the real element in
        Section~\ref{sec:single-element-validation}.
\end{itemize}

A standard half-wave dipole presents a radiation resistance near $Z_d\approx
\SI{73}{\ohm}$ in isolation, but this value drops substantially once the
element is loaded by closely spaced parasitic directors, often falling well
below \SI{30}{\ohm}. Such a low, geometry-sensitive impedance is difficult to
match robustly over a useful band. The receiving element of this project
therefore uses a folded dipole as its driven element, a choice that is
standard in commercial UHF television antennas for exactly this reason
\citep{kraus1988}.

Because the radiating current of a folded dipole is shared between its two
parallel conductors, the current at the feed terminals for the same radiated
power is halved compared to a simple dipole, and, since impedance is voltage
over current, halving the current for the same power raises the impedance by
a factor of four. This gives the approximate relation
\begin{equation}
    Z_f \;\approx\; 4\,Z_d \;\approx\; 4 \times 73\,\si{\ohm}
    \;\approx\; 292\,\si{\ohm}.
    \label{eq:folded-dipole}
\end{equation}
Equation~\eqref{eq:folded-dipole} \citep{balanis2016,kraus1988} is only an
approximation: it holds for two parallel conductors of equal radius, with a
small spacing relative to a wavelength, in a balanced configuration, and
\emph{before} the strong loading introduced by a reflector and multiple
directors is taken into account. It is a design-level estimate of the folded
dipole \emph{on its own}, not the final impedance measured at the antenna
port once the complete Yagi structure is assembled; that final value, which
does include the director and reflector loading, is $Z_\mathrm{in}$, reported
in Section~\ref{sec:single-element-validation}.

Two practical advantages follow from the estimate of Eq.~\eqref{eq:folded-dipole}
nonetheless. First, the resulting impedance ($\sim\!\SI{300}{\ohm}$) is high
enough that the parasitic loading of the directors brings it down into a
convenient range rather than collapsing it toward zero, giving a more stable,
broadband match. Second, the $\sim\!\SI{300}{\ohm}$ balanced impedance is
transformed to the standard \SI{75}{\ohm} unbalanced coaxial feed by a simple
$4{:}1$ balun, which simultaneously suppresses common-mode currents on the
outer conductor of the cable. The combination of a folded dipole, a $4{:}1$
balun, and a \SI{75}{\ohm} feed is the canonical UHF-television front end,
and it is inherited unchanged by the present instrument. This is also why
\SI{75}{\ohm} is used as the reference impedance $Z_0$ throughout this
chapter, in the reflection-coefficient and SWR relations of
Eq.~\eqref{eq:swr}: the antenna is a \SI{75}{\ohm} DVB-T2 device, the
mast-head LNA is a \SI{75}{\ohm} part, and the entire feed is \SI{75}{\ohm}
coaxial cable, so \SI{75}{\ohm} is the impedance against which return loss
and SWR are physically meaningful for this system, both in simulation
(Section~\ref{sec:simulation}) and in the NanoVNA measurement
(Section~\ref{sec:single-element-validation}).

% =====================================================================
\section{Reference Observing System: Learmonth/RSTN and the Xuan~Mai Field Setup}
\label{sec:reference-systems}

The \SI{610}{\mega\hertz} flux record used throughout this thesis
(Chapter~\ref{ch:theory}, Section~\ref{subsec:sv-of-t}) is not produced by the
antenna characterised in this chapter; it comes from a separate, professional
observing system at Learmonth. Before describing the field hardware built for
this project, it is worth being precise about what that reference system
actually is and, just as importantly, what its public documentation does and
does not specify.

% ---------------------------------------------------------------------
\subsection{The Learmonth/RSTN Reference System}
\label{subsec:learmonth-hardware}

Learmonth Solar Observatory, in Western Australia, is one of the stations of
the Radio Solar Telescope Network (RSTN), a network of ground-based solar
radio patrol stations originally established by the U.S. Air Force Research
Laboratory in the 1960s, with official solar patrols beginning in
1966~\citep{rstn_overview}. Each RSTN station produces two distinct kinds of
data: discrete-frequency total-power measurements from a
\emph{Radio Interference Monitoring Set} (RIMS), and wideband spectral data
from a separate \emph{Solar Radio Spectrograph} (SRS) covering the metre and
decametre range~\citep{rstn_overview}. The Learmonth RIMS instrument monitors
the Sun continuously on eight discrete frequencies-245, 410, 610, 1415, 2695,
4995, 8800, and \SI{15400}{\mega\hertz}-using three parabolic dish antennas
of different sizes: the largest, \SI{8.5}{\metre} (28~ft) in diameter, feeds
the \SI{245}{}, \SI{410}{}, and \SI{610}{\mega\hertz} channels used in this
thesis; a \SI{2.4}{\metre} (8~ft) dish covers 1415-\SI{8800}{\mega\hertz}; and
a \SI{1}{\metre} (3~ft) dish covers \SI{15400}{\mega\hertz}~\citep{sws_rims_metadata}.
Each channel's radiometer output is digitised at a \SI{1}{\hertz} sampling
rate and transferred to the Sydney World Data Centre for processing and
archival~\citep{sws_rims_metadata}. The RIMS one-second data record, at each
monitored frequency, "the total power output in SFU\ldots the total raw power
received by the RIMS"~\citep{rstn_overview}: it is a calibrated, single-channel
total-power measurement, not a raw voltage or IQ stream, and daily Noon Flux
values are additionally tagged with a quality qualifier (good; weather
attenuated; interference; unknown; burst) that flags questionable
data~\citep{rstn_overview}.

This description supports the way Learmonth data are used in this thesis, but
it also marks the limit of what can be claimed. The publicly available
Learmonth/RSTN documentation establishes the antenna, the frequency plan, the
sampling rate, and the fact that the reported quantity is total power in
solar flux units; it does not publish the detailed analogue front-end or IF
filter response of the \SI{610}{\mega\hertz} channel. Accordingly, this
thesis makes no claim about the exact receiver bandwidth or filter shape used
at Learmonth: the public metadata identifies the channel as a
discrete-frequency RIMS measurement, but does not specify the detailed RF/IF
filter response~\citep{sws_rims_metadata,rstn_overview}. The archived
$S_\nu(t)$ record is therefore used, correctly, as an external, already-calibrated
flux input (Section~\ref{subsec:sv-of-t}), and no attempt is made in this
thesis to reproduce the Learmonth receiver chain itself.

% ---------------------------------------------------------------------
\subsection{Preliminary Field Deployment at Xuan~Mai}
\label{subsec:field}

To test whether solar radio observation was feasible with this hardware, a
preliminary two-element interferometer was deployed at the Citrus Research and
Development Centre in Xuan~Mai, Chuong~My, Hanoi, over several sessions in
December~2025. Two identical HKD~G11KD-T2 antennas were mounted on heavy-duty
photographic tripods and separated along an approximately east-west baseline of
$\sim\!\SI{8}{\metre}$ (Figure~\ref{fig:xuan-mai-field}). The tripods allowed
manual azimuth and elevation adjustment to follow the solar transit. Each
antenna was fitted with a mast-head Low-Noise Amplifier (LNA) powered through a
custom dual-channel bias-tee, which injected DC voltage up the coaxial cable
while passing the \SI{610}{\mega\hertz} RF signal to the receiver. The back-end
electronics consisted of an SDRplay~RSPdx-R2 software-defined radio receiver
controlled by a Raspberry~Pi~4 single-board computer. All components were housed
in a weatherproof enclosure with a \SI{12.6}{\volt} battery pack for autonomous
field operation (Figure~\ref{fig:receiver-node}).

\begin{figure}[H]
  \centering
  \includegraphics[width=\textwidth]{photo_xuan_mai_field}
  \caption[Two-element interferometer field deployment at Xuan~Mai (December 2025)]{Two-element interferometer field deployment at Xuan~Mai (December 2025).\\[4pt]
  Two HKD~G11KD-T2 antennas mounted on heavy-duty photographic tripods and
  separated along an approximately east-west baseline of $\sim\!\SI{8}{\metre}$
  at the Citrus Research and Development Centre, Chuong~My, Hanoi. Each element
  is equipped with a mast-head LNA powered through a custom bias-tee; the
  \SI{610}{\mega\hertz} signal is routed via coaxial cable to the
  SDRplay~RSPdx-R2 receiver at the processing node. Manual azimuth and
  elevation adjustment allowed the beams to track the solar transit.}
  \label{fig:xuan-mai-field}
\end{figure}

\begin{figure}[H]
  \centering
  \includegraphics[width=\textwidth]{photo_receiver_node}
  \caption[Receiver and processing node: interior view of the field deployment enclosure]{Receiver and processing node: interior view of the field deployment enclosure.\\[4pt]
  Top-down view of the weatherproof enclosure housing all back-end electronics.
  The SDRplay~RSPdx-R2 software-defined radio receiver (centre) digitises the
  \SI{610}{\mega\hertz} signal arriving via the \textbf{RF~in} port through a
  dual-channel bias-tee that simultaneously injects \SI{12}{\volt} DC up the
  coaxial feed to power the mast-head LNAs. A Raspberry~Pi~4 single-board
  computer runs the data-acquisition software and writes recorded IQ streams to
  a \SI{1}{\tera\byte} HDD. A \SI{12.6}{\volt} lithium battery pack with a
  5\,V buck converter provides autonomous field power; a \SI{12}{\volt} fan
  maintains thermal margin inside the sealed enclosure.}
  \label{fig:receiver-node}
\end{figure}

The campaign confirmed that the instrument concept was sound in principle-the
data showed a broad diurnal rise consistent with the Sun transiting the antenna
beams-but it also exposed a series of operational failures that prevented any
reliable detection of a genuine solar radio burst. Table~\ref{tab:learmonth-vs-xuanmai}
summarises how this local, uncalibrated test setup compares with the
professional Learmonth/RSTN reference system just described.

\begin{table}[H]
  \centering
  \caption[Learmonth/RSTN reference system versus the Xuan~Mai field setup]{Learmonth/RSTN reference system versus the Xuan~Mai field setup.}
  \label{tab:learmonth-vs-xuanmai}
  \renewcommand{\arraystretch}{1.3}
  \begin{tabularx}{\textwidth}{l >{\raggedright\arraybackslash}X >{\raggedright\arraybackslash}X}
    \toprule
    \textbf{Aspect} & \textbf{Learmonth/RSTN (reference)} & \textbf{Xuan~Mai (this thesis)} \\
    \midrule
    Antenna & \SI{8.5}{\metre} parabolic dish (245/410/\SI{610}{\mega\hertz} group) & Two HKD~G11KD-T2 Yagi-Uda elements, $\sim\!\SI{9.5}{dBi}$ each \\
    Receiver & RIMS radiometer, calibrated, professionally maintained & SDRplay~RSPdx-R2 SDR with mast-head LNA, consumer-grade \\
    Reported quantity & Total power, calibrated to solar flux units, \SI{1}{\hertz} & Raw complex (IQ) voltage stream, uncalibrated \\
    Timing/phase reference & Not applicable (total-power radiometer, not an interferometer) & Independent receiver oscillators, no shared clock \\
    Data quality control & IFLUX qualifier flags (good / weather / interference / unknown / burst) & None automated; RFI, cable loss, and hardware faults diagnosed post-hoc (Section~\ref{subsec:element-vs-system}) \\
    Role in this thesis & Source of the real flux input $S_\nu(t)$ (Chapter~\ref{ch:theory}) & Motivates the need for a controlled antenna model (this chapter) \\
    \bottomrule
  \end{tabularx}
\end{table}

The relationship between the three systems introduced so far can be summarised
in one sentence: \emph{Learmonth validates the solar input; Xuan~Mai motivates
the hardware model; the 4NEC2 antenna simulator bridges the two by providing a
controlled antenna response and coupling estimate.}

% ---------------------------------------------------------------------
\subsection{Element Versus System: Diagnosing the Field Failure}
\label{subsec:element-vs-system}

The most important outcome of the field campaign, for the purposes of this
thesis, is a clean separation between the performance of the receiving
\emph{element} and the performance of the receiving \emph{system}. The element
was not at fault; the system was. The evidence is as follows.

\paragraph{The element is validated.} During troubleshooting, the antennas and their
LNAs were characterised directly at the feed point with a vector network
analyser (NanoVNA). Each element returned a return loss
$S_{11}\approx\SI{-12.8}{dB}$ and a voltage standing-wave ratio
$\mathrm{VSWR}\approx1.59$ at \SI{610}{\mega\hertz}, and the LNAs delivered their
rated gain. These figures indicate a well-matched, correctly functioning
receiving element. The radiating hardware, in other words, met its specification-answering
question~1 of Section~\ref{sec:why-model} in the affirmative.

\paragraph{The system was degraded at the physical layer.} The failures lay in
the chain between the antenna terminals and the digitiser:
\begin{enumerate}[leftmargin=1.8em]
    \item \textbf{Cable insertion loss.} Transmission-coefficient ($S_{21}$)
    measurements of the coaxial runs alone showed
    \SIrange{-3.1}{-3.5}{dB} of loss. More than half of the received power was
    therefore dissipated in the feed before reaching the receiver-fatal for a
    source already close to the noise floor.
    \item \textbf{Bias-tee / LNA fault.} One mast-head LNA lost power through a
    disconnection in its bias-tee channel, leaving that arm of the
    interferometer effectively dead during part of the campaign.
    \item \textbf{Terrestrial RFI and ground noise.} The outdoor noise floor was
    densely populated with carriers from nearby electronic equipment and
    broadcast services. As anticipated in Section~\ref{subsec:suitability}, the
    DVB-T2 antenna's in-band sensitivity made this interference unavoidable, and
    the recorded dynamic spectra showed strong, non-solar fringe-like structure
    and night-time oscillations driven by local sources.
    \item \textbf{Short, ill-defined baseline and weak source.} The achievable
    physical baseline ($\sim\!\SI{8}{\metre}$ in the field) and the faintness of
    the quiet-Sun signal ($\sim\!0.5~\si{dB}$ excursions above the receiver noise floor) left little margin
    against the degradations above.
    \item \textbf{Absence of phase coherence between receiver chains.} The
    two RTL-SDR dongles each had an independent \SI{28.8}{\mega\hertz} crystal
    oscillator with no shared reference. This caused a random, time-varying
    phase error $\phi_{\rm err}(t)$ between the two arms. When the
    cross-product $V_1 \cdot V_2^*$ was averaged over time, this random phase
    swept across $[0,2\pi]$ and drove the fringe amplitude toward zero — a
    phenomenon known as \emph{fringe washing} \citep{tms2017}. Without a
    shared reference clock or phase calibration on a known source, fringes
    cannot be recovered no matter how strong the solar signal is. This is the
    most fundamental barrier to interferometry with independent receivers.
\end{enumerate}

\paragraph{Results:} The element is confirmed to work correctly, while the
system was degraded by interference, cable loss, and — as discussed in
Chapter~\ref{ch:results} — the lack of phase coherence between two independent
receiver chains. Given this, the right approach is not to keep rebuilding the
field array by trial and error. Instead, the antenna is characterised on its own
in a controlled numerical environment where the geometry is known to the
millimetre and where cabling, interference, and fabrication tolerances do not
apply. A rigorous electromagnetic model provides the true primary beam and the
coupling values that the field experiment could not supply directly. It also
provides the clean, idealised receiving model on which the rest of the
simulation depends. The remainder of this chapter builds and validates that
model.

% =====================================================================
\section{4NEC2 Antenna Simulation Methodology}
\label{sec:simulation}

To obtain precise electromagnetic parameters free of the field's confounding
variables, the U-frame folded-dipole Yagi was modelled in 4NEC2: a free antenna
simulation software package that provides a graphical implementation of the
Numerical Electromagnetics Code (NEC), based on the Method of Moments
\citep{nec2}. In the rest of this chapter it is referred to, interchangeably,
as the 4NEC2 antenna simulator or the 4NEC2 Method-of-Moments antenna model.

% ---------------------------------------------------------------------
\subsection{The Method of Moments}
\label{subsec:mom}

The 4NEC2 antenna simulator solves the thin-wire Electric-Field Integral
Equation (EFIE) by the Method of Moments (MoM) \citep{harrington1993,nec2}. The
continuous antenna is discretised into short, straight wire segments; the unknown
current distribution is expanded in basis functions on these segments; and
enforcement of the boundary condition that the total tangential electric field
vanish on the conducting surface converts the integral equation into a dense
linear system $[\,Z\,]\,[\,I\,] = [\,V\,]$, where $[\,Z\,]$ is the generalised
impedance matrix, $[\,V\,]$ the excitation, and $[\,I\,]$ the segment currents
to be solved. Accurate results require each segment to be short compared with
the wavelength, typically between $\lambda/10$ and $\lambda/20$, and the
segment radius to satisfy the thin-wire approximation. Once the current
distribution is known, the input impedance follows from the feed-segment voltage
and current, and the far-field pattern from the radiation integral over all
segments. This NEC wire-model simulation is exceptionally well suited to open
wire structures such as a Yagi-Uda array: it discretises only the conductors
(not the surrounding volume), converges rapidly, and yields impedance and
pattern data of high accuracy in the UHF band.

% ---------------------------------------------------------------------
\subsection{Wire Model Construction}
\label{subsec:model}

The HKD geometry of Table~\ref{tab:hkd-geom} was translated into a
fourteen-wire NEC model for the 4NEC2 antenna simulator: the four U-frame
reflector bars, the folded-dipole driven element (modelled as its constituent
wires), and the six directors, with conductor radii in the
\SIrange{2.0}{3.5}{\milli\metre} range to match the physical tubing
(Figure~\ref{fig:4nec2-model}). The model was solved in free space (ground flag
$\texttt{GN}\,{-}1$); this is the physically correct choice for a
sky-pointing radio-astronomy antenna, whose intrinsic pattern is the quantity
required, with the effect of the real ground treated separately rather than
baked into the element pattern. The excitation was a \SI{1}{\volt} source
placed at the centre of the folded dipole (the central segment of the
driven-element wire), i.e.\ at the physical feed point, and, as established
in Section~\ref{sec:matching}, the impedance reference throughout the model
was set to \SI{75}{\ohm}, matching the antenna, LNA, and coaxial feed used in
the field.

Table~\ref{tab:hkd-geom} gives the physical dimensions in engineering terms;
Figure~\ref{fig:4nec2-geometry-table} shows the same information as it is
actually entered into the 4NEC2 antenna simulator, wire by wire, before the
model is solved.

\begin{figure}[H]
  \centering
  \includegraphics[width=\textwidth]{fig_4nec2_geometry_table}
  \caption[4NEC2 Geometry editor: wire-by-wire coordinate input for the HKD~G11KD-T2 model]{4NEC2 Geometry editor: wire-by-wire coordinate input for the HKD~G11KD-T2 model.\\[4pt]
  Screenshot of the 4NEC2 antenna simulator's Geometry tab for the fourteen-wire
  model: the four reflector bars (wires~1-4), the two folded-dipole conductors
  and their two short end caps (wires~5-8), and the six directors
  (wires~9-14). Each row gives the wire's segment count (\textbf{Segs}, the
  Method-of-Moments discretisation of Section~\ref{subsec:mom}), its two
  endpoint coordinates $(X_1,Y_1,Z_1)$ and $(X_2,Y_2,Z_2)$ in metres, and its
  conductor radius. The short end-cap wires (7 and~8, three segments each)
  close the folded-dipole loop; every other wire uses 21 segments, short
  enough compared with $\lambda\approx\SI{491.5}{\milli\metre}$ to satisfy the
  thin-wire, short-segment requirement of the Method of Moments. This table
  is the direct numeric realisation of the physical dimensions surveyed in
  Table~\ref{tab:hkd-geom}; the rendered result of solving this exact
  geometry is shown in Figure~\ref{fig:4nec2-model} below.}
  \label{fig:4nec2-geometry-table}
\end{figure}

\begin{figure}[H]
  \centering
  \includegraphics[width=\textwidth]{fig_4nec2_wiremodel}
  \caption[4NEC2 Method-of-Moments wire model of the HKD~G11KD-T2 antenna]{4NEC2 Method-of-Moments wire model of the HKD~G11KD-T2 antenna.\\[4pt]
  The fourteen-wire NEC geometry represents the four U-frame reflector bars
  (rear group), the two-wire folded-dipole driven element at the feed, and the
  six directors extending toward the front. Wire radii range from
  \SIrange{2.0}{3.5}{\milli\metre} to match the physical tubing dimensions of
  Table~\ref{tab:hkd-geom}. The model is solved in free space
  (\texttt{GN\,{-}1}) with a \SI{1}{\volt} excitation placed at the central
  segment of the driven element.}
  \label{fig:4nec2-model}
\end{figure}

% =====================================================================
\section{Single-Element Validation: Impedance, Gain, Beamwidth, SWR}
\label{sec:single-element-validation}

Solving the model at \SI{610}{\mega\hertz} yields the single-element parameters
collected in Table~\ref{tab:ant-params-sim}. This section answers questions~2
and~3 of Section~\ref{sec:why-model}: what the real beam looks like, and
whether the match is stable enough for the received power to be trusted.

\begin{table}[H]
  \centering
  \caption[Simulated single-element parameters at 610~MHz (4NEC2, free space)]{Simulated single-element parameters at 610~MHz (4NEC2, free space).\\[4pt]
           The eleven-element U-frame Yagi-Uda antenna (HKD G11KD-T2) was
           solved in free space by the Method of Moments at the operating
           frequency. All impedance and SWR values are referenced to the
           75~$\Omega$ coaxial feed standard; gain and beamwidth are evaluated
           in the E-plane.}
  \label{tab:ant-params-sim}
  \renewcommand{\arraystretch}{1.25}
  \begin{tabularx}{\textwidth}{l >{\raggedright\arraybackslash}X}
    \toprule
    \textbf{Parameter} & \textbf{Value} \\
    \midrule
    Operating frequency               & \SI{610}{\mega\hertz} ($\lambda\approx\SI{491.5}{\milli\metre}$) \\
    Input impedance $Z_\mathrm{in}$   & $108.4 - j22.0~\si{\ohm}$ \\
    Reflection coeff. $|\Gamma|$ (at \SI{75}{\ohm}) & $\approx 0.215$ ($\approx\SI{-13.3}{dB}$) \\
    SWR (at \SI{75}{\ohm})            & $1.55$ \\
    Forward gain                      & \SI{9.48}{dBi} \\
    Half-power beamwidth (E-plane)    & $\approx\SI{52}{\degree}$ \\
    Radiation efficiency              & $\approx\SI{100}{\percent}$ (lossless, free space) \\
    \bottomrule
  \end{tabularx}
\end{table}

The input impedance is $Z_\mathrm{in}=(108.4 - j22.0)\,\si{\ohm}$: a real part close
to the \SI{75}{\ohm} reference and a small capacitive reactance indicating the
antenna is operating just below its exact resonance. As anticipated in
Section~\ref{sec:matching}, this is the impedance of the complete Yagi
structure, not the isolated folded-dipole estimate $Z_f$ of
Eq.~\eqref{eq:folded-dipole}: director and reflector loading have pulled it
well below the $\sim\!\SI{300}{\ohm}$ folded-dipole value. Referenced to
\SI{75}{\ohm}, the reflection-coefficient magnitude is
\begin{equation}
    |\Gamma| = \left|\frac{Z_\mathrm{in}-Z_0}{Z_\mathrm{in}+Z_0}\right|
             = \left|\frac{(108.4-j22.0)-75}{(108.4-j22.0)+75}\right|
             \approx 0.215,
    \label{eq:gamma-single}
\end{equation}
equivalent to a return loss of about \SI{-13.3}{dB}. Substituting $|\Gamma|=0.215$
into Eq.~\eqref{eq:swr} gives $\mathrm{SWR}=(1+0.215)/(1-0.215)\approx 1.55$.
An SWR of $1.55$ (comfortably below the conventional $2.0$ acceptability
threshold) means that the great majority of incident power is delivered into the
feed with negligible reflective loss, and that no additional matching network is
required-the folded-dipole/$4{:}1$-balun front end matches the array to
\SI{75}{\ohm} on its own. This directly answers question~3 of
Section~\ref{sec:why-model}: the received power computed from the Friis
relation (Eq.~\eqref{eq:friis}) is not being silently discounted by an
unmatched feed.

The antenna delivers a forward gain of \SI{9.48}{dBi}, nearly an order of
magnitude above an isotropic radiator, with an E-plane half-power beamwidth of
about \SI{52}{\degree} (Figure~\ref{fig:pattern}). The full three-dimensional
directivity pattern (Figure~\ref{fig:3d-pattern}) confirms that this gain is
concentrated into a single dominant end-fire main lobe with the rear suppressed
by the U-frame reflector, and that no significant sidelobe structure redirects
sensitivity toward the ground or the horizon. This beamwidth is a good
compromise for an un-tracked or hand-tracked solar instrument: narrow enough to
reject horizon noise and to give the U-frame reflector its high front-to-back
ratio, yet wide enough that the Sun dwells within the main lobe for a useful
fraction of the day without continuous mechanical steering.

\begin{figure}[H]
  \centering
  \includegraphics[width=\textwidth]{fig_pattern_eplane}
  \caption[Simulated E-plane radiation pattern of the Yagi element (4NEC2, free space)]{Simulated E-plane radiation pattern of the Yagi element (4NEC2, free space).\\[4pt]
           The end-fire main lobe directed toward the Sun is clearly visible,
           along with the high front-to-back ratio achieved by the U-frame
           reflector, which suppresses ground emission entering from behind.
           The half-power beamwidth of approximately $\SI{52}{\degree}$ is wide
           enough for the Sun to remain within the main lobe for several hours
           without mechanical steering, yet narrow enough to reject horizon
           noise and ground spillover.}
  \label{fig:pattern}
\end{figure}

\begin{figure}[H]
  \centering
  \includegraphics[width=\textwidth]{fig_3d_radiation_pattern}
  \caption[Three-dimensional directivity pattern of the HKD~G11KD-T2 at \SI{610}{\mega\hertz} (4NEC2)]{Three-dimensional directivity pattern of the HKD~G11KD-T2 at \SI{610}{\mega\hertz} (4NEC2 antenna simulator).\\[4pt]
  The colour scale runs from $-41$~dBi (deep blue) to $+9.48$~dBi (red),
  confirming the peak forward gain reported in Table~\ref{tab:ant-params-sim}.
  The dominant end-fire main lobe (red--yellow) accounts for the majority of
  the radiated power and is directed toward the sky during solar transit.
  The smaller rear lobe (blue) is substantially suppressed by the U-frame
  reflector, yielding a high front-to-back ratio that limits ground thermal
  noise pickup. The green arrows indicate the antenna boom axis orientation.}
  \label{fig:3d-pattern}
\end{figure}

% ---------------------------------------------------------------------
\subsection{Cross-Validation: Measurement, MoM, and FEM}
\label{subsec:validation}

A simulated result is only reliable if it is confirmed by other means. A
particular strength of this study is that the single-element impedance match is
confirmed by \emph{three independent methods} based on completely different
principles: a direct physical measurement using a NanoVNA vector network
analyser at the feed point (Figure~\ref{fig:nanovna-setup}), the
Method-of-Moments simulation in the 4NEC2 antenna simulator, and an independent
finite-element result from previous work on the same antenna.
Table~\ref{tab:validation} shows all three.

\begin{figure}[H]
  \centering
  \includegraphics[width=\textwidth]{photo_nanovna_setup}
  \caption[NanoVNA measurement at the HKD~G11KD-T2 feed point (Xuan~Mai)]{NanoVNA measurement at the HKD~G11KD-T2 feed point (Xuan~Mai).\\[4pt]
  The NanoVNA vector network analyser is connected directly to the antenna
  feed terminals via a short SMA-to-coaxial adapter, with the antenna oriented
  skyward to reproduce operational conditions. Port~1 of the VNA drives the
  driven element and records the reflected signal, from which $S_{11}$ and the
  \SI{75}{\ohm}-referenced SWR at \SI{610}{\mega\hertz} are extracted.
  This measurement constitutes the ``NanoVNA (measured)'' row of
  Table~\ref{tab:validation}.}
  \label{fig:nanovna-setup}
\end{figure}

\begin{table}[H]
  \centering
  \caption[Three-way SWR cross-validation at 610~MHz]{Three-way SWR cross-validation at 610~MHz.\\[4pt]
           The standing-wave ratio referenced to 75~$\Omega$ is compared across
           three independent methods: a direct VNA measurement at the Xuan~Mai
           deployment site, the 4NEC2 Method-of-Moments solution obtained in
           this work, and a Finite Element Method result from previous work on
           the same physical antenna. All three results fall within the narrow
           range 1.22--1.59, confirming that the simulated parameters are
           reliable.}
  \label{tab:validation}
  \renewcommand{\arraystretch}{1.25}
  \begin{tabularx}{\textwidth}{l >{\raggedright\arraybackslash}X c c}
    \toprule
    \textbf{Method} & \textbf{Principle} & \textbf{SWR (\SI{75}{\ohm})}
                    & \textbf{$S_{11}$} \\
    \midrule
    NanoVNA (measured)    & physical (VNA)        & $1.59$        & $\approx\SI{-12.8}{dB}$ \\
    4NEC2 (this work)     & Method of Moments     & $1.55$        & $\approx\SI{-13.3}{dB}$ \\
    HFSS (previous work)  & Finite Element Method & $1.22$--$1.53$ & $\approx\SI{-13.6}{dB}$ to $\SI{-20}{dB}$ \\
    \bottomrule
  \end{tabularx}
\end{table}

All three results fall within the narrow SWR range of $1.2$--$1.6$, with the
MoM result from this work sitting between the physical measurement and the
finite-element solution. Agreement between a direct measurement and two solvers
based on different numerical methods (integral-equation MoM versus volumetric
FEM) is strong evidence that the extracted parameters are reliable.

It is worth discussing the small differences between the three results rather
than presenting an unrealistically perfect match. The measured SWR ($1.59$) is
slightly worse than the simulated value ($1.55$). This is expected: the
physical measurement includes balun insertion loss, connector mismatches, small
fabrication tolerances, and the effect of the real measurement environment —
none of which exist in the idealised free-space model. The FEM result covers a
slightly wider range because it was obtained over several geometry iterations.
Despite these differences, all three converge to within a few percent of each
other in SWR, which justifies using the simulated pattern in the imaging
pipeline.

% =====================================================================
\section{Mutual-Coupling Analysis and Impedance-Matrix Interpretation}
\label{sec:coupling-rationale}

Section~\ref{sec:single-element-validation} characterised one antenna on its
own. This section asks the fourth and last question of
Section~\ref{sec:why-model}: when a second, identical antenna is placed
nearby to form an interferometric baseline, does it disturb the first one
enough to matter?

% ---------------------------------------------------------------------
\subsection{From a Single Antenna to a Coupled Pair}
\label{subsec:coupling-intuition}

When an antenna stands entirely on its own, the voltage at its feed depends
only on its own current, through its self-impedance:
\begin{equation}
    V_1 = Z_{11} I_1 .
    \label{eq:v1-alone}
\end{equation}
Now bring a second, identical antenna into the picture, separated from the
first by some baseline distance. Antenna~2 carries its own current $I_2$,
and-exactly as for the parasitic elements of a single Yagi in
Eq.~\eqref{eq:mutual-matrix}-that current radiates a field that induces an
additional voltage on antenna~1. The feed voltage of antenna~1 is now
\begin{equation}
    V_1 = Z_{11} I_1 + Z_{12} I_2 ,
    \label{eq:v1-coupled}
\end{equation}
and, by the same argument applied to antenna~2,
\begin{equation}
    V_2 = Z_{21} I_1 + Z_{22} I_2 .
    \label{eq:v2-coupled}
\end{equation}
Written together as a two-port impedance matrix,
\begin{equation}
    \begin{bmatrix} V_1 \\ V_2 \end{bmatrix}
    =
    \begin{bmatrix} Z_{11} & Z_{12} \\ Z_{21} & Z_{22} \end{bmatrix}
    \begin{bmatrix} I_1 \\ I_2 \end{bmatrix},
    \label{eq:two-port-matrix}
\end{equation}
where:
\begin{itemize}[leftmargin=1.5em,itemsep=1pt,topsep=2pt]
  \item $Z_{11}$ and $Z_{22}$ are the \emph{self-impedances} of antennas~1 and~2,
        equal to the isolated $Z_\mathrm{in}$ of Section~\ref{sec:single-element-validation}
        when the two antennas are identical;
  \item $Z_{12}=Z_{21}$ (for identical, reciprocal antennas) is the
        \emph{mutual impedance}: the voltage induced on one antenna's feed
        per unit current on the other.
\end{itemize}
If $Z_{12}$ is small compared with $Z_{11}$, the two antennas behave almost
independently, exactly as required by the imaging theory of
Chapter~\ref{ch:theory}. If $Z_{12}$ is not small, the beam, the impedance
match, and-most importantly for an interferometer-the measured phase of each
antenna are all disturbed by the presence of the other.

The impedance actually seen at antenna~1's feed, once both antennas are
simultaneously excited, is the \emph{active} (or driving-point) impedance,
\begin{equation}
    Z_{\mathrm{active},1} = \frac{V_1}{I_1} = Z_{11} + Z_{12}\,\frac{I_2}{I_1} ,
    \label{eq:driving-impedance}
\end{equation}
which depends on the current ratio $I_2/I_1$ and therefore on how the pair is
driven. Two special cases isolate $Z_{11}$ and $Z_{12}$ individually, and they
are also the two cases that can each be run directly as a single simulation
in the 4NEC2 antenna simulator:
\begin{itemize}[leftmargin=1.5em,itemsep=1pt,topsep=2pt]
  \item \textbf{Even mode} ($I_2=I_1$, both ports driven in phase):
        $Z_\mathrm{even}=Z_{11}+Z_{12}$;
  \item \textbf{Odd mode} ($I_2=-I_1$, ports driven \SI{180}{\degree} out of
        phase): $Z_\mathrm{odd}=Z_{11}-Z_{12}$.
\end{itemize}
Solving this pair of equations gives $Z_{11}$ and $Z_{12}$ separately from two
independent 4NEC2 runs:
\begin{equation}
    Z_{11} = \tfrac{1}{2}\big(Z_\mathrm{even}+Z_\mathrm{odd}\big),
    \qquad
    Z_{12} = \tfrac{1}{2}\big(Z_\mathrm{even}-Z_\mathrm{odd}\big).
    \label{eq:z11-z12-solve}
\end{equation}
This even/odd technique, applied in Section~\ref{subsec:coupling-results}
below, is the reason two simulations-not one-are run at every element
separation.

% ---------------------------------------------------------------------
\subsection{The Independent-Sensor Assumption in Aperture Synthesis}
\label{subsec:independence}

The imaging framework adopted in this thesis-the van~Cittert-Zernike theorem
and the dirty-image machinery of Chapter~\ref{ch:theory}-relies on a specific
assumption: that each receiving element samples the spatial coherence
function of the sky \emph{independently}, with a primary beam $A_n(l,m)$
identical to the isolated single-element pattern \citep{tms2017}. In the
language of Section~\ref{subsec:coupling-intuition}, this is precisely the
statement that $Z_{12}$ is negligible compared with $Z_{11}$ at the array's
operating baselines: when it holds, the measured cross-correlation between
two elements is a clean sample of the complex visibility, and the recovered
phase is purely the geometric delay between the two apertures, with no
additional contribution from the antennas disturbing one another.

% ---------------------------------------------------------------------
\subsection{Why Coupling Specifically Corrupts Interferometric Phase}
\label{subsec:coupling-phase}

For a total-power instrument, the main consequence of coupling would be a
mismatch loss. For an \emph{interferometer} the consequence is more insidious,
because the measurement is fundamentally a measurement of phase. Coupling
perturbs the active impedance of Eq.~\eqref{eq:driving-impedance} (mismatch
and standing waves), distorts and makes asymmetric the primary beam, and-most
damagingly-adds a parasitic, geometry-dependent phase term to the received
signal. The interferometer requires that the measured phase be the pure
geometric delay between the elements; any additional coupling-induced phase
contaminates the visibility and propagates directly into the dirty image as
spurious structure \citep{tms2017,gupta1983}.

The standard ways to quantify coupling are equivalent and well established
\citep{balanis2016,pozar2012}: the off-diagonal mutual impedance $Z_{12}$; the
transmission scattering parameter $S_{21}$ (the same information expressed in
incident/reflected waves, and the quantity most directly measurable with a
network analyser); the full $N\!\times\!N$ coupling matrix for an
$N$-element array; and the shift of the active impedance relative to the
isolated element. It is therefore not sufficient to \emph{assume} the elements
are independent at the operating baseline-the assumption must be
\emph{demonstrated} by computing one of these measures as a function of
element separation. That demonstration follows next.

% ---------------------------------------------------------------------
\subsection{Array Layout and Mutual-Coupling Extraction}
\label{subsec:coupling-results}

To quantify the inter-element coupling, a two-port model was built in the
4NEC2 antenna simulator by duplicating the validated single element and
translating the copy along the baseline (east-west) axis, leaving the
individual antenna geometry untouched. Four separations were studied:
\SI{0.5}{\metre}, \SI{5}{\metre}, \SI{10}{\metre}, and \SI{15}{\metre}.

Both ports were excited with a \SI{1}{\volt} source in phase (the
\emph{even mode} of Section~\ref{subsec:coupling-intuition}). Under even-mode
excitation each port presents the active impedance
$Z_\mathrm{active}=Z_{11}+Z_{12}$, so the shift of $Z_\mathrm{active}$ away
from the isolated value is read directly, in a single run, as the fingerprint
of the mutual impedance $Z_{12}$. The even-mode run additionally yields the
two-element array pattern. To obtain the magnitude of the transmission
coefficient $|S_{21}|$ precisely-which cannot be read from a single run-the
even/odd technique of Eq.~\eqref{eq:z11-z12-solve} was used, combining the
in-phase run with a \SI{180}{\degree} out-of-phase (odd-mode) run. The results
are given in Table~\ref{tab:coupling}.

\begin{table}[H]
  \centering
  \caption[Mutual coupling versus east-west element separation (4NEC2)]{Mutual coupling versus east-west element separation (4NEC2 antenna simulator).\\[4pt]
           Two U-frame Yagi elements were simulated at four separations using
           the even/odd excitation technique. $|\Delta Z|$ is the
           active-impedance shift magnitude relative to the isolated value
           ($Z_\mathrm{in}=108.4-j22.0\,\si{\ohm}$). Coupling falls from
           $-34$~dB at 0.5~m ($\approx1\lambda$) to below $-78$~dB at 5~m
           and below $-90$~dB at 15~m element separation, formally
           establishing the independent-sensor assumption.}
  \label{tab:coupling}
  \renewcommand{\arraystretch}{1.3}
  \begin{tabular}{l c c c c}
    \toprule
    \textbf{Baseline} & \textbf{$d/\lambda$} &
    \textbf{$Z_\mathrm{in}$, ant.~1 ($\si{\ohm}$)} &
    \textbf{$|\Delta Z|$ ($\si{\ohm}$)} & \textbf{$|S_{21}|$ (dB)} \\
    \midrule
    isolated        & -   & $108.4 - j22.0$ & $0$     & -  \\
    \SI{0.5}{\metre} & $1.0$  & $108.0 - j18.0$ & $3.94$  & $-34$ \\
    \SI{5}{\metre}   & $10.2$ & $108.4 - j21.9$ & $0.028$ & $-78$ \\
    \SI{10}{\metre}  & $20.3$ & $108.4 - j22.0$ & $0.010$ & $-87$ \\
    \SI{15}{\metre}  & $30.5$ & $108.4 - j22.0$ & $0.007$ & $-90$ \\
    \bottomrule
  \end{tabular}
\end{table}

The behaviour is unambiguous. At \SI{0.5}{\metre} ($\approx\!1\lambda$) the two
elements still share near-field structure: the active reactance shifts from
$-22.0$ to $-18.0~\si{\ohm}$, a $3.94~\si{\ohm}$ ($3.6\%$) change, and the
isolation is $|S_{21}|\approx\SI{-34}{dB}$. This is the deliberately
``worst case'' point, included to show the coupling regime where it is genuinely
present. Once the separation reaches \SI{5}{\metre} ($\approx\!10\lambda$),
however, the coupling has already collapsed: the impedance shift falls below
$0.03~\si{\ohm}$ ($<0.03\%$) and the isolation drops past
\SI{-78}{dB}. By \SI{15}{\metre} the impedance perturbation is within numerical
noise and $|S_{21}|\approx\SI{-90}{dB}$.

\begin{figure}[H]
  \centering
  \includegraphics[width=\textwidth]{fig_coupling_curve}
  \caption[Mutual-coupling isolation and impedance shift versus element separation]{Mutual-coupling isolation and impedance shift versus element separation.\\[4pt]
           The isolation $|S_{21}|$ and active-impedance shift $|\Delta Z|$
           are both plotted against east-west element separation. Both quantities
           collapse rapidly beyond $1\lambda$: by 5~m ($\approx10\lambda$) the
           isolation exceeds 78~dB and $|\Delta Z|<0.03\,\si{\ohm}$, confirming
           that elements at the operational baselines behave as fully independent
           sensors.}
  \label{fig:coupling}
\end{figure}

The operational baselines of the simulated interferometer start at
\SI{8}{\metre}; the imaging pipeline used in the remainder of this thesis is
configured for baselines of \SI{8}{\metre}, \SI{16}{\metre}, and \SI{24}{\metre} ($\approx\!48.8\lambda$), which lies
far deeper still in the decoupled regime. Extrapolating the monotonic decay of
Table~\ref{tab:coupling}, the isolation at \SI{24}{\metre} is well below
\SI{-90}{dB} and the active-impedance perturbation is entirely negligible.
Across the full range of operating baselines, therefore, the elements neither
perturb each other's impedance nor distort each other's radiation pattern at any
level that could affect the measurement, and the parasitic coupling phase
discussed in Section~\ref{subsec:coupling-phase} is absent. The
independent-sensor assumption underlying the van~Cittert-Zernike imaging is thus
not merely assumed but quantitatively established-question~4 of
Section~\ref{sec:why-model} is answered-and the signal-processing
pipeline may treat each antenna as an ideal, independent sensor whose measured
phase is pure geometric delay.

% =====================================================================
\section{Primary-Beam Extraction for Chapter~\ref{ch:pipeline}}
\label{sec:primary-beam-extraction}

The final product of the electromagnetic model is the primary beam consumed by
the imaging pipeline. The single-element far-field pattern was exported from
the 4NEC2 antenna simulator's output over a $181\times361$ grid in
$(\theta,\phi)$ (a $1^{\circ}$ step in both elevation and azimuth) and
normalised to unity at boresight. This normalised pattern is the angular
response multiplied onto each antenna's received voltage in the pipeline, so
that the simulated sky is sampled through the true, physically derived
antenna response rather than through an idealised isotropic or Gaussian beam-directly
answering question~2 of Section~\ref{sec:why-model}.

A point of method deserves emphasis. The pattern exported is that of the
\emph{single element} (\SI{9.48}{dBi}), \emph{not} that of the two-element array
(\SI{12.49}{dBi}) produced in the even-mode coupling runs. This is deliberate
and physically necessary: in an interferometer the array factor is synthesised
by the \emph{cross-correlation} of the two independent element signals, not by
the radiation pattern of a combined aperture. Embedding the two-element array
pattern as the primary beam would double-count the inter-element geometry-once
in the beam and again in the correlation-and corrupt the visibilities. The
correct primary beam is therefore the isolated element pattern, exactly as the
negligible-coupling result of Section~\ref{subsec:coupling-results} entitles us
to use. The exported file (\texttt{UFRAME\_primary\_beam\_610MHz.npz}) is loaded
directly by the pipeline of Chapter~\ref{ch:pipeline}: it is read by the
\texttt{PrimaryBeam} class and applied at every simulation time step to scale
the received signal amplitude via Eq.~\eqref{eq:psig}. Every downstream
result-transit power, visibility amplitude, dirty-image brightness-therefore
reflects the actual Yagi beam characterised in this chapter.

With this, all four questions posed in Section~\ref{sec:why-model} have been
answered by direct evidence rather than by assumption: the element receives a
well-matched \SI{610}{\mega\hertz} signal (Section~\ref{sec:single-element-validation});
its real, simulated beam-not a Gaussian placeholder-is what the pipeline uses
(this section); its impedance match is stable and threefold cross-validated
(Section~\ref{subsec:validation}); and mutual coupling at every operational
baseline is negligible, so the independent-sensor assumption of
Chapter~\ref{ch:theory} holds (Section~\ref{sec:coupling-rationale}). Recalling
the summary of Section~\ref{sec:reference-systems}, Learmonth supplied the
validated solar input, the Xuan~Mai field campaign supplied the motivation, and
the 4NEC2 antenna simulator supplied the controlled antenna response and
coupling estimate that together make the pipeline of the next chapter
trustworthy.
\clearpage
